\section{Discussion}
\label{sec:discussion}

\subsection{Contrast with the tetramer chain}

The main contrast with the reference tetramer follows from unit-cell parity.
Four spin-$1/2$ degrees of freedom can form an integer-spin or singlet cluster sector.
Five leave a half-integer degree of freedom unless translation symmetry is broken or the ground state is degenerate.
At the tested positive-$\alpha$ pentamer points, the surviving doublet converts the inter-cell exchange into an antiferromagnetic interaction between effective spins.
The finite-size data are therefore consistent with a critical effective chain, whereas tetramer phase names and its local triplon, quarton, or higher composite labels cannot be transferred by geometric analogy alone~\cite{li2026tetramer_verified}.

The spectral comparison sharpens this distinction.
The pentamer's odd-parity excited doublet is dark at $q=\pi$, while higher even-parity spin-$3/2$ sectors remain bright.
In the chain, that local hierarchy coexists with a dispersive low-energy continuum and five-site folding or high-energy structure.
The multiplet analysis predicts visibility, while the many-body calculation determines whether the weight forms a continuum, a broadened branch, or a more complex folded response.

\subsection{Controlled interpretation near \texorpdfstring{$\alpha=0$}{alpha=0}}

For $|\alpha|\ll\min(1,\beta)$, a trimer--dimer decomposition gives a useful limiting picture.
Projecting the trimer doublets and eliminating the inter-cell dimer triplet produces the candidate exchange
\begin{equation}
\Jeff=\frac{2\alpha^2}{9\beta}>0.
\label{eq:near-zero}
\end{equation}
An eight-spin trimer--dimer--trimer oracle evaluated at 47 points supports this local $\alpha^2/\beta$ scaling in its asymptotic domain, consistent with the effective-spin logic of verified trimer-chain work~\cite{cheng2022trimer}.
The finite-$\alpha$ controls at $(-0.10,0.25)$ and $(-0.175,0.25)$ deviate from the leading term by $-36.1351\%$ and $-53.8550\%$, respectively.
These are signed deviations of a local controlled approximation, not numerical errors in the $L=80$ calculation.
The open-boundary sector gap also need not equal a bulk excitation gap.
The decomposition therefore provides a mechanism close to $\alpha=0$, not proof of a second phase or a global $c=1$ region.

\subsection{Meaning and boundary of the negative-\texorpdfstring{$\alpha$}{alpha} result}

The negative-$\alpha$ scan answers a narrower but reproducible question: the dominant finite-size momentum changes across the coupling plane.
It does not determine whether that change sharpens into a Lifshitz transition, belongs to a broad incommensurate regime, or remains a finite-size crossover.
The local quartet analysis supplies a candidate origin for the deep-region value: the dominant folded momentum of the isolated spin-$3/2$ cell is $0.60\pi$ throughout the negative region, matching the measured $\qpeak$ there, while the $\alpha>0$ doublet projects onto $q=\pi$ instead.
That correspondence is suggestive rather than conclusive, because a local shape factor fixes a momentum scale without determining whether the extended chain develops a singularity at it, and because the near-zero strip shows the largest sensitivity to initialization and fitting window.
It is, however, a falsifiable statement of the kind the present data set could not resolve: a controlled larger-size calculation of the deep-region momentum should converge to the local value if the reconstruction is set by the single-cell multiplet.
The decisive limitation of the phase question is not a failed DMRG convergence test.
The accepted Stage C cells pass the numerical and provenance gates, and the most sensitive endpoint passes the $k=8\rightarrow12$ solver control.
The limitation is physical: the momentum and finite-size-flow diagnostics do not identify the same boundary interval.

Numerical convergence validates the state and measured observables at a finite size; it cannot manufacture agreement between physically distinct diagnostics.
Conversely, the absence of a joint boundary does not prove that the entire negative region is a single phase.
Our bounded conclusion is that the finite-size momentum reconstruction is established, while a separately named thermodynamic phase is not.

\subsection{Remaining discriminating tests}

Two limitations bound the present work and define the calculations that would move it forward.
First, no bond-dimension series was performed.
The Krylov controls fix an eigensolver setting, not the matrix-product-state bond dimension $\chi$, so the entanglement and gap observables are validated against solver noise but not against truncation of the state space itself.
A $\chi$ series at the two most sensitive points, $(\alpha,\beta)=(-0.10,0.25)$ and $(-0.175,0.25)$ at $L=80$, is the single most informative missing calculation.
Second, the historical negative-$\alpha$ size sequence mixes solver configurations, as noted in Sec.~\ref{sec:methods}; recomputing the $L=40,60$ endpoints with the current setting would be required before any $\ceff$ or gap series is given a finite-size interpretation.

A stronger phase claim would additionally require an observable with a controlled large-size singularity or crossing that is consistent with the momentum reconstruction.
A negative-$\alpha$ dynamical representative could test whether the momentum reorganization is accompanied by a qualitative spectral transfer, but such a calculation would strengthen the phenomenology rather than prove a boundary by itself.
